% D.88 reserved candidate. Notes only; D.87 approval and full audits pending.
\subsection{Every finite relative local depth in a zero-five-valuation cusp packet}

\begin{statement}{D.88}{Open}{Candidate: the compact kernel at every finite cusp index}
Let \(a\ne0\) be an arbitrary completed Kummer class over
\(\mathbb Q\) with zero valuation at five. Write
\[
 r=\min_{q\ne5}v_5(a_q),\qquad
 \lambda=\sum_{q\ne5}a_q\frac{\log(q^4)}{4\log6},
 \qquad N=v_5(\lambda)<\infty,
 \qquad e=N-r\ge0.                                    \tag{D.88.1}
\]
The valuation support may be infinite \(c_0\) support. For the
actual cusp pair \(s=(a,0)\), the compact groups of D.63 are
\[
 \mathcal S_s\simeq
 \begin{cases}
  0,&e=0,\\
  \displaystyle\prod_{j\ge0}\mathbb Z/5^e,&e>0,
 \end{cases}
 \qquad
 \mathcal K_s\simeq\mathcal S_s\times\mathbb Z_5^{2e}.
                                                               \tag{D.88.2}
\]
These are compact topological isomorphisms; the splitting is
continuous and noncanonical. More precisely D.77 identifies
\(L_5\to\mathcal T_s\) with the identity on the zero-parameter
countable product and, on the other two boundaries, the two
projections \(\mathbb Z_5^N\to\mathbb Z_5^r\).

For every \(r\ge0,e\ge1\), the actual rational tangential
cusp parameter
\[
                    a=5^r\kappa(1+5^{e+1})             \tag{D.88.3}
\]
has precisely these indices. Its nonempty metabelian fibre is
a torsor under \(\prod_{\mathbb N}\mathbb Z_5\), by D.74.
Thus every bounded five-primary exponent in (D.88.2) occurs
in a genuine cusp section with finite ramification.

The statement requires \(a_5=0\) and \(N<\infty\).
It does not handle the additional actual valuation-graph
relation when \(a_5\ne0\), or infer the infinite endpoint
from equal infinite coranks. It makes no claim about the first
higher-derived obstruction or all noncuspidal abelian pairs.
\end{statement}

\begin{proof}[Candidate proof, pending full independent audits]
The new input D.87 identifies the finite-global-localization
kernel with the \emph{entire} actual operator-image restriction
kernel. We will compute that kernel, retaining its finite
torsion and all actual coefficient actions.

\paragraph{Exact images and the relative cyclotomic module.}
Use the same notation \(\Lambda,t_1,t_2,u,v,w,P,I_0,\Gamma\)
as D.86. For \(n\ge0\), put
\[
                     q_n=(1+u)^{5^n}-1,
                     \qquad q_0=u.
\]
The actual global and local images, after the integral local
gauge of D.83, are
\[
 I_s=5^r\mathbb Z_5\rtimes I_0,
       \qquad I_5=5^N\mathbb Z_5\rtimes I_0.           \tag{D.88.4}
\]
The common complement has the untwisted action \(\rho_0\),
because \(a_5=0\). If \(e=0\), the two actual images agree
and D.83, or D.87, gives \(\mathcal S_s=0\). D.77 then
also gives \(\mathcal K_s=0\). Assume \(e>0\).

Set
\[
 B=\Lambda/(q_r),\quad R=\Lambda/(q_N),\quad J=q_rR,
 \qquad F_{r,N}=q_N/q_r,
 \qquad A=\Lambda/(F_{r,N}).                           \tag{D.88.5}
\]
Multiplication by the \emph{actual} generator \(q_r\)
identifies the scalar module underlying \(J\) with \(A\).
The polynomial \(F_{r,N}\) is monic, has constant term
\(5^e\), and is independent of \(v\). Thus \(A\) is
finite free over \(\mathbb Z_5[[v]]\), and
\[
                     A/(u)=(\mathbb Z/5^e)[[v]].       \tag{D.88.6}
\]
Multiplication by \(u\) is injective in \(A\). Indeed,
if \(uf=F_{r,N}g\) in the full power-series ring, evaluation
at \(u=0\) gives \(5^e g(0,v)=0\), so \(g=ug_1\);
cancellation in the full domain gives \(f=F_{r,N}g_1\).
This is valid even though \(A\) has several cyclotomic
components and its quotient (D.88.6) has five-torsion.

Continuous coinduction and D.87 give the natural formula
\[
 \mathcal S_s^*\simeq
       (J^D)^{I_0}/\operatorname{im}((R^D)^{I_0}).       \tag{D.88.7}
\]
On the cyclotomic kernel the primitive cubic element of D.77
annihilates precisely \(PJ\), and all other beta differences
are multiples of \(P\). Consequently
\[
 (J^D)^{I_0}=(T(-1)_\Gamma)^*,\qquad
             T=J/PJ\simeq A/(uvw)                    \tag{D.88.8}
\]
in scalar coordinates. The retained actual generator is
\(q_r\). In particular its substitution factor is
\(S_\chi(t_1^{5^r})\), with constant coefficient \(\chi\),
not a scalar obtained by dividing by \(5^r\).

\paragraph{The exact torsion layer; the snake map is zero.}
As \(v\) and \(w\) form a regular sequence in \(A\),
there is an exact sequence of compact actual modules
\[
 0\longrightarrow E\longrightarrow T\longrightarrow Q
      \longrightarrow0,
 \quad E=vwA/(uvwA)\simeq(\mathbb Z/5^e)[[v]],
 \quad Q=A/(vw).                                      \tag{D.88.9}
\]
For the regular-sequence assertion, \(v\) is a nonzero
 divisor in the formal variable and \(w\bmod v=u\) is
injective by the previous paragraph. The actual generator of
\(E\) is \(q_rvw\). At \(u=0\), the Alexander multiplier
contributes \(\chi S_\chi(t_2)\), substitution on \(q_r\)
contributes \(\chi\), and substitution on each of \(v,w\)
contributes \(S_\chi(t_2)\). After the one Cartier untwist,
the action on (D.88.6) is
\[
 f(v)\longmapsto\chi S_\chi(1+v)^3
                         f((1+v)^\chi-1).             \tag{D.88.10}
\]
Multiplication by \(v^3\) identifies this with
\(v^3C\), where \(C=(\mathbb Z/5^e)[[\mathbb Z_5]]\)
has action \(t^x\mapsto\chi t^{\chi x}\).
The missing first three degrees have nontrivial tame
characters \(\chi,\chi^2,\chi^3\). Exact order-four
averaging therefore gives
\[
                     E(-1)_\Gamma\simeq
                           \prod_{j\ge0}\mathbb Z/5^e.
                                                               \tag{D.88.11}
\]
Indeed its invariant discrete dual is the space of continuous
functions \(f:\mathbb Z_5\to\mathbb Z/5^e\) satisfying
\(f(\gamma x)=\gamma^{-1}f(x)\). Such a function vanishes
at zero, has one freely chosen value on each nonzero valuation
shell, and has only finitely many nonzero shell values by
continuity at zero. This proves the asserted compact topology.

There is no hidden finite loss in (D.88.11). On
\[
             Q/uQ=(\mathbb Z/5^e)[[v]]/(v^2)
\]
the actual Cartier-untwisted action has successive characters
\(\chi^2,\chi^3\): the retained generator is \(q_r\).
Their order-four invariant parts are zero. Hence
\(Q^\Gamma\subset uQ\). The natural map
\[
                              uT\longrightarrow uQ
                                                               \tag{D.88.12}
\]
is an equivariant isomorphism. Surjectivity is immediate;
its kernel is \(uT\cap E\), which is zero because modulo
\(u\) the product \(vw\) is \(v^2\), a nonzero divisor
even over \(\mathbb Z/5^e\). Equivalently,
\((u)\cap(vw)=(uvw)\) inside \(A\).
Thus every invariant element of \(Q\) has an invariant
lift in \(uT\).

Choose an actual lift of \(2\in\Gamma\), which topologically
generates \(\Gamma\). Apply the snake lemma to its operator
minus one in (D.88.9). The preceding invariant lifting makes
the connecting map zero, so there is an exact sequence
\[
 0\longrightarrow E(-1)_\Gamma\longrightarrow T(-1)_\Gamma
       \longrightarrow Q(-1)_\Gamma\longrightarrow0. \tag{D.88.13}
\]
This proves injectivity of the infinite torsion layer; it does
not discard a finite snake-map image by comparing ranks.

\paragraph{The actual finite boundary and its torsion.}
Let \(R_n=\mathbb Z_5[\mathbb Z/5^n]\), including
\(R_0=\mathbb Z_5\), and let
\[
       J_{r,N}=\ker(R_N\longrightarrow R_r)=q_rR_N.
\]
The action is the genuine permutation action
\(t\mapsto t^\chi\). Restriction to the two boundaries
\(v=0,w=0\) gives an injection of \(Q\) into their two
scalar \(A\)-modules. Its cokernel is
\(A/(v,w)=\mathbb Z/5^e\), with actual untwisted
character \(\chi^2\), so its tame invariant and coinvariant
parts vanish.

On the first boundary, multiplication by \(u q_r\)
identifies the actual Cartier-untwisted scalar module with
\(uJ_{r,N}\subset R_N\). On the diagonal boundary one
also multiplies by \(t^{-1/2}\), using D.39's exact
reflection formula. These identifications are integral:
if \(u q_r f\in q_N\Lambda\), cancellation of \(q_r\),
evaluation at \(u=0\), and cancellation of \(u\) show
\(f\in F_{r,N}\Lambda\). They preserve the exact action
because substitution supplies the factors
\(S_\chi(t)S_\chi(t^{5^r})\).

The quotient
\(J_{r,N}/uJ_{r,N}\simeq\mathbb Z/5^e(1)\)
has nontrivial tame character. Exact tame averaging gives
\[
       (uJ_{r,N})_\Gamma\simeq(J_{r,N})_\Gamma.        \tag{D.88.14}
\]
We compute the latter integrally. The \(\Gamma\)-orbits in
\(\mathbb Z/5^n\) are zero and its \(n\) nonzero valuation
shells. Its group-ring invariants are the orbit sums, and its
coinvariants are free on the orbit classes. On invariants,
the map \(R_N\to R_r\) sends the zero basis vector to
itself and sends the shell of valuation \(j<r\) to
\(5^e\) times the corresponding shell sum. Shells with
\(j\ge r\) map to a multiple of the zero basis vector.
Its invariant cokernel is therefore
\( (\mathbb Z/5^e)^r\).
On coinvariants the same map is onto: it maps the first
\(r\) shells to the corresponding shells and all remaining
ones to zero's orbit class. Its kernel is free of rank
\(N-r=e\). The exact sequence for the operator \(2-1\)
thus gives
\[
 0\longrightarrow(\mathbb Z/5^e)^r
   \longrightarrow(J_{r,N})_\Gamma
   \longrightarrow\mathbb Z_5^e\longrightarrow0.
\]
The free quotient splits continuously, yielding
\[
 Q(-1)_\Gamma\simeq
                (\mathbb Z/5^e)^{2r}\times\mathbb Z_5^{2e}.
                                                               \tag{D.88.15}
\]
The finite summand is essential when \(r>0\); it has not
been removed by a rational dimension calculation.

\paragraph{The whole restriction kernel and its exponent.}
Equations (D.88.11), (D.88.13), and (D.88.15) imply
\[
 T(-1)_\Gamma\simeq T_{\rm tor}\times\mathbb Z_5^{2e},
\]
where \(T_{\rm tor}\) is a compact bounded-exponent module
and fits naturally in an exact sequence
\[
 0\longrightarrow\prod_{j\ge0}\mathbb Z/5^e
    \longrightarrow T_{\rm tor}
    \longrightarrow(\mathbb Z/5^e)^{2r}\longrightarrow0.
                                                               \tag{D.88.16}
\]
For example boundedness by \(5^{2e}\) follows immediately
from this sequence. The free quotient is split by choosing
lifts of its finite basis.

By D.77's natural map, the quotient
\((R^D)^{I_0}/(B^D)^{I_0}\) is the divisible group
\((\mathbb Q_5/\mathbb Z_5)^{2e}\). Its injective image
in \((J^D)^{I_0}\) is the whole maximal divisible subgroup:
that subgroup has the same finite corank \(2e\), and
ordinary duality turns the injection into a surjection between
free \(\mathbb Z_5\)-modules of the same rank. Thus
(D.88.7) identifies \(\mathcal S_s\) with \(T_{\rm tor}\).

In fact this whole module is killed by \(5^e\), not merely
by \(5^{2e}\). The actual subgroup \(I_5\subset I_s\)
is open of index \(5^e\). Restriction followed by
corestriction on degree-one continuous cohomology is
multiplication by that index. Here is the cocycle verification
with no normality assumption. For left coset representatives
\(t\), write \(gt=t_g h(g,t)\). If \(z\) is a global
cocycle, the corestriction of its restriction has value
\[
 \sum_t t_g z(h(g,t))
   =5^e z(g)+(g-1)\sum_t z(t),
                                                               \tag{D.88.17}
\]
using the cocycle identity and the permutation \(t\mapsto t_g\).
The same finite-sum formula is a continuous cocycle and sends
coboundaries to coboundaries. Therefore every class in the
restriction kernel is killed by \(5^e\). By D.87 this is
the entire \(\mathcal S_s^*\), hence also
\(5^e\mathcal S_s=0\).

The exact sequence (D.88.16) is now a sequence of modules over
\(\mathbb Z/5^e\), whose quotient is finite free over that
ring. Choosing lifts of its basis splits it continuously.
Its finite free summand is absorbed by reindexing the countable
product. We obtain
\[
                       \mathcal S_s\simeq
                            \prod_{j\ge0}\mathbb Z/5^e.
\]
This is the asserted exact compact group, including its exponent.
The argument does not silently identify a quotient by an
arbitrary finite subgroup with the same bounded product.

\paragraph{The free local factor and actual examples.}
The two nonzero boundary classes \(a,a-b\) have global
depth \(r\) and local depth \(N\); the zero class \(b\)
contributes the identity on its countable product. D.77 therefore
gives the stated natural map, its kernel \(\mathbb Z_5^{2e}\),
and the continuous splitting in (D.88.2).

For (D.88.3), the integer \(1+5^{e+1}\) is congruent to two
modulo four, so its valuation at two is one. It is a unit at
five and its principal-unit logarithm has valuation \(e+1\).
Since \(v_5(\log6)=1\), its Kummer class is globally primitive
and has local unit coefficient of valuation \(e\).
Multiplication by \(5^r\) gives exactly (D.88.1).
The cusp inertia construction supplies an actual full section
with this parameter, and finite support permits application
of D.74 to its nonempty metabelian fibre.
\end{proof}
